In this section, we describe the empirical setting and research design, explain the construction of the analysis sample, and assess the validity of the randomization and compliance with treatment assignment.

\subsection{Research Design}

We study the effect of digital information provision on household electricity consumption using a framed field experiment \citep{listWhyEconomistsShould2011} conducted in cooperation with the largest energy utility in Upper Austria. Households were provided with smart meter data at 15-minute intervals and, depending on the experimental condition, received access to a smartphone application offering detailed feedback on electricity use. The application was designed to make otherwise opaque consumption patterns more visible and easier to interpret by translating raw electricity use into indicators such as consumption levels, monetary costs, and carbon emissions.

Recruitment took place in three phases between May and November 2017. Approximately 30,000 email invitations and 10,000 postal letters were distributed, each including a description of the app and a link to an online registration survey.\footnote{Postal invitations were discontinued after the first recruitment phase because of a low response rate of below 2\%.} Of 1,939 initial respondents, 1,524 completed the registration process. After excluding ineligible households, 1,020 households in the control and information arms were randomly assigned within tranche using a random number generator. Participants provided sociodemographic information, dwelling characteristics, and appliance ownership data and consented to the use of their smart meter data.\footnote{The study was approved by the Austrian Data Protection Authority (reference number DSB-D036.500/0005-DSB/2017).}

The present paper focuses on the comparison between the information condition and the control group. The broader field trial also included additional treatment arms that combined app access with dynamic tariff incentives. Because those arms bundle information with monetary incentives, they are not informative about the information-only effect that is the focus of this paper and are therefore excluded from the present analysis. Our main estimand is thus the intention-to-treat effect of being assigned access to the application.

Treated households received access to the app in June (tranche 1), September (tranche 2), and November 2017 (tranche 3), while control households remained untreated throughout the study period. Table~\ref{tab:recruitment} summarizes recruitment and the final analysis sample.

\begin{table}[H]
\centering
\caption{Recruitment and final analysis sample by tranche}
\label{tab:recruitment}
\begin{tabular}{llcccccc}
\toprule
 &  & \multicolumn{3}{c}{Initial sample} & \multicolumn{3}{c}{Final analysis sample} \\
\cmidrule(lr){3-5} \cmidrule(lr){6-8}
Tranche & Start date & Control & Information & Total & Control & Information & Total \\
\midrule
1 & 06/06/2017 & 166 & 200 & 366 & 154 & 171 & 325 \\
2 & 19/09/2017 & 224 & 224 & 448 & 185 & 196 & 381 \\
3 & 20/11/2017 & 116 & 90  & 206 & 104 & 82  & 186 \\
\midrule
Total &  & 506 & 514 & 1020 & 443 & 449 & 892 \\
\bottomrule
\end{tabular}

\caption*{\footnotesize
Note: The table reports the number of households recruited and included in the final analysis sample by tranche and experimental condition. Differences between the initial and final samples arise from data cleaning steps applied after randomization, including the exclusion of households with photovoltaic systems, insufficient data coverage, anomalous consumption patterns, and extreme outliers.}
\end{table}

\subsection{Treatment}

Households were randomly assigned to one of two conditions in the present analysis.\footnote{Randomization was implemented by the utility under the supervision of the academic project partners using a random number generator.}

\begin{itemize}
    \item \textbf{Information condition ($N = 449$):} households received access to a smartphone application providing detailed feedback on electricity consumption.
    \item \textbf{Control condition ($N = 443$):} households did not receive access to the application.
\end{itemize}

The application was available through standard app stores and provided households with feedback on electricity consumption in kilowatt-hours (kWh), corresponding monetary costs in euros, and carbon emissions in kilograms of \(CO_2\). These indicators were based on approximate conversion factors. Due to the data transmission process, information was typically available with a delay of up to 24 hours, which reflects the standard timing of smart meter data availability in Austria. Users could review their consumption across different time horizons (day, week, month, year), facilitating the interpretation of usage patterns and making electricity use more visible in everyday decision-making.

\begin{figure}[H]
    \centering
    \includegraphics[width=0.7\linewidth]{paper/figures/MILP_app.png}
    \caption{Overview of the app interface displaying consumption, costs, and emissions.}
    \label{fig:app_overview}
\end{figure}

Beyond individual feedback, the application included a social comparison feature that allowed households to benchmark their consumption against similar households. 

Expected consumption was estimated with a simple OLS model on observable household characteristics. Each household was shown the expected value of its own consumption conditional on those characteristics, i.e. a benchmark against statistically similar peers rather than the full sample.

The app also incorporated gamification and learning elements intended to promote engagement. Households could predict their electricity consumption for the following day and received feedback on the accuracy of their predictions. In addition, the app provided information on the electricity use of specific appliances and activities. These features did not involve monetary rewards, but included a points-based system designed to encourage repeated interaction.

\begin{figure}[H]
    \centering
    \includegraphics[width=0.55\linewidth]{paper/figures/MILP_app_game.png}
    \caption{Gamification feature allowing households to predict and learn about electricity consumption.}
    \label{fig:app_gamification}
\end{figure}

\subsection{Data and Analysis Sample}

We measure electricity consumption using smart meter data recorded at 15-minute intervals for all households in the sample. These high-frequency measurements allow us to observe detailed consumption patterns over time and capture short-term behavioral responses to the intervention. For the empirical analysis, we aggregate the 15-minute observations to the hourly or weekly level.

The final analysis sample was constructed through several data-cleaning steps applied after randomization. First, 109 households with photovoltaic (PV) systems were excluded because reliable information on self-consumption and electricity generation was unavailable. 

Second, we restricted the sample to households with sufficiently complete meter data, requiring at least 90\% coverage of potential observations. The roughly 10\% of observations that are missing stem from technical failures in the smart-meter and data-transmission infrastructure rather than from household behaviour.

Third, days with an unusually high number of zero-consumption readings, which were typically caused by area-wide communication failures between smart meters and the data platform rather than genuine zero usage, were removed.

Fourth, we excluded households with extreme average electricity consumption, defined as more than three standard deviations above the mean in the pre-treatment period. Finally, to ensure comparability across groups within tranche, we imposed a common pre-treatment window. The resulting analysis sample contains 892 households.

Because treatment was assigned before these restrictions were applied, it is important to assess whether sample selection is correlated with treatment status. We therefore estimate a series of logistic regressions of exclusion indicators on treatment assignment. Across all cleaning steps, we find no evidence that exclusion of households differs between conditions. A detailed breakdown of the sample construction and the corresponding selection diagnostics is reported in Appendix Table~\ref{tab:sample_selection}. Moreover, due to the staggered rollout of the intervention, the length of the pre-treatment period differs across tranches. The pre-treatment window comprises 20, 7, and 12 days of observations in tranches 1, 2, and 3, respectively, while the end of the observation period is common across tranches.

\subsection{Randomization checks}

Before assessing the impact of information provision on household electricity consumption, we conduct several checks to evaluate the validity of the empirical design. Table \ref{table:balance_check} shows that the treatment and control groups are highly comparable across key household characteristics. Home ownership is common in both groups (75\% in the control group and 73\% in the information group), and average dwelling size is similar at 134.4 and 128.9 square meters, respectively. With respect to heating technologies, around one-third of households have a heat pump (33\% vs.\ 30\%), while electric heating is relatively rare (6\% vs.\ 5\%). The prevalence of energy-intensive appliances is also balanced across groups: 24\% versus 21\% of households own a swimming pool, 19\% versus 17\% report a sauna, and air conditioning is uncommon in both groups (4\% vs.\ 3\%). Electric vehicles are rare as well, with only around 2\% of households reporting ownership in either group. None of these differences is statistically significant at conventional levels, which supports the validity of the random assignment.  This conclusion is reinforced by a regression of treatment assignment on household characteristics and tranche fixed effects. The covariates are jointly insignificant (Wald test, \(p = 0.61\)), providing no indication that treatment assignment is systematically related to observed household characteristics.


\begin{table}[H]
\centering
\caption{Covariate Balance Between Control and Information Groups}
\label{table:balance_check}
\begin{threeparttable}
\small
\begin{tabular}{lccc}
\toprule
 & Control & Information & $p$-value \\
\midrule

Households & 443 & 449 &  \\

\midrule
\textbf{Heating} \\
Gas              & 0.23 (0.42) & 0.23 (0.42) & 0.899 \\
District heating & 0.14 (0.35) & 0.16 (0.37) & 0.346 \\
Biomass          & 0.11 (0.31) & 0.12 (0.33) & 0.577 \\
Oil              & 0.17 (0.38) & 0.15 (0.35) & 0.362 \\
Electric heating & 0.06 (0.23) & 0.05 (0.21) & 0.515 \\
Heat pump        & 0.33 (0.47) & 0.30 (0.46) & 0.433 \\

\midrule
\textbf{Housing} \\
Home ownership      & 0.75 (0.43) & 0.73 (0.44) & 0.424 \\
Apartment           & 0.31 (0.46) & 0.28 (0.45) & 0.388 \\
Single-family house & 0.05 (0.23) & 0.07 (0.25) & 0.357 \\
Semi-detached house & 0.62 (0.49) & 0.63 (0.48) & 0.666 \\

\midrule
\textbf{Appliances} \\
Swimming pool   & 0.24 (0.43) & 0.21 (0.41) & 0.287 \\
Sauna           & 0.19 (0.39) & 0.17 (0.38) & 0.430 \\
Air conditioning& 0.04 (0.19) & 0.03 (0.17) & 0.547 \\
Aquarium        & 0.06 (0.24) & 0.05 (0.21) & 0.426 \\
Dryer           & 0.57 (0.50) & 0.58 (0.49) & 0.757 \\
Water bed       & 0.05 (0.22) & 0.04 (0.20) & 0.490 \\
Deep freezers   & 0.82 (0.69) & 0.83 (0.71) & 0.884 \\
Computers       & 2.50 (1.40) & 2.57 (1.40) & 0.492 \\
Electric vehicle& 0.02 (0.14) & 0.02 (0.15) & 0.836 \\
E-bike          & 0.07 (0.33) & 0.04 (0.23) & 0.118 \\

\midrule
\textbf{Other} \\
Number of residents       & 2.67 (1.20)   & 2.69 (1.20)   & 0.761 \\
Dwelling size (m$^2$)     & 134.42 (65.36)& 128.92 (54.14)& 0.171 \\

\bottomrule
\end{tabular}

\begin{tablenotes}[flushleft]
\footnotesize
\item The table reports means by experimental group, with standard deviations in parentheses. Binary variables are reported as shares. $p$-values are based on two-sided tests of equality in means between the control and information groups.
\end{tablenotes}

\end{threeparttable}
\end{table}

We next examine pre-treatment consumption patterns. Figure~\ref{fig:pre_treatment_load_curve} plots average hourly electricity consumption by treatment status during the pre-treatment period. The load profiles of the two groups are closely aligned throughout the day. To account for staggered rollout, Figure~\ref{fig:tranche_pre_calendar} presents pre-treatment consumption separately for each tranche in calendar time. Across all tranches, the control and information groups exhibit similar pre-treatment patterns. These graphical patterns are also reflected in the summary statistics. Mean hourly consumption is highly similar across groups within tranche (for example, 0.359 versus 0.370 kWh in tranche 1 and 0.413 versus 0.412 kWh in tranche 2), and differences in total pre-treatment consumption are small relative to within-group variation. Taken together, the balance checks and pre-treatment comparisons support the internal validity of the experimental design.

\begin{figure}[H]
\centering
\includegraphics[width=1\linewidth]{paper/figures/descriptive_loadcurve_hourly_pre.pdf}
\caption{Average hourly electricity consumption in the pre-treatment period by treatment group.}
\label{fig:pre_treatment_load_curve}
\end{figure}

\subsection{Compliance}

Compliance with treatment assignment was incomplete because treated households had to install and actively use the application. Among households in the information condition in the final analysis sample, 65\% accessed the app at least once. Interaction rates were 47\% for the social comparison feature and 40\% for the gamification feature. The main information page was the most frequently used component, followed by gamification and social comparison. App engagement was highly skewed and persistent only for a subset of users. The top 10\% of households accounted for 76.7\% of all recorded interactions, while the median household interacted with the app only six times. Moreover, 34.5\% of treated households never used the app after installation, and an additional 43\% engaged only briefly, for between one and three months. Only 22.7\% of households exhibited sustained engagement over four or more months. Figure~\ref{fig:engagement} shows the evolution of app usage over the post-treatment period. These compliance patterns imply that the treatment estimates reported in the paper should be interpreted as intention-to-treat effects of being assigned access to the app rather than as causal effects of actual app use. 

A post-trial survey of treated households (\(N = 111\)) provides additional descriptive evidence on the perceived usefulness of the different app features. Respondents rated the information page as the most useful component, followed by social comparison and gamification. These responses should, however, be interpreted cautiously because survey participation was voluntary and non-random.

\begin{figure}[H]
    \centering
    \includegraphics[width=1\linewidth]{paper/figures/app_engagement.pdf}
    \caption{App engagement over time during the post-treatment period.}
    \label{fig:engagement}
\end{figure}